\chapter{Tám Chuyện Cuối Bàn}
\chapterintro{Cuối lớp thường là nơi câu chuyện phát triển nhanh hơn bài giảng}{Nhiều cuộc nói chuyện bắt đầu bằng một câu nhỏ rồi dài đúng một tiết.}

\begin{lucbatbox}{Lục bát: Cuối bàn mở hội}
Cuối bàn như có giờ chơi\
Một câu khều nhẹ kéo mười câu sang\
Thầy cô giảng ở trên bảng\
Dưới này chuyện đã sang trang thứ mười\
Nếu mà chia nửa tiếng cười\
Cho phần bài học, hẳn đời bớt run
\end{lucbatbox}

\begin{tutuyetbox}{Thất ngôn tứ tuyệt: Nguồn tin cuối lớp}
Cuối bàn cập nhật rất nhanh nhé
Tin nào mới nổi hóa đề chung
Chỉ riêng bài học trong vùng ấy
Hay đi lạc giữa tiếng cười thôi
\end{tutuyetbox}

\begin{ghichu}
Cuối bàn không xấu. Chỉ là nếu không giữ nhịp, nó rất dễ biến thành trung tâm phát sóng chuyện riêng của cả lớp.
\end{ghichu}

\ngancach

\begin{lucbatbox}{Lục bát: Cười khúc khích rất đều}
Một câu chưa nói thật to\
Cả bàn đã khúc khích cho đủ đầy\
Bạn trên quay xuống chau mày\
Thầy cô dừng lại một giây nhìn về\
Hóa ra niềm vui quá nghề\
Cũng làm tiết học nặng nề ít nhiều
\end{lucbatbox}

\begin{tutuyetbox}{Thất ngôn tứ tuyệt: Trò riêng quá đà}
Chuyện riêng em kể rất mê say
Say hơn phần giảng hôm nay nhiều
Nếu đem nhiệt ấy vào chuyên cần
Cuối bàn hẳn đã sáng điểm rồi
\end{tutuyetbox}